\thispagestyle{plain}

\chapter{Introduction}
\label{cp:introduction}

\vspace{2 \baselineskip}
\section{Objective of the Project}
The primary goal of this project is to identify key factors associated with income levels using the 1994 \emph{Adult} dataset. We aim to build and evaluate a binary classification model that predicts whether an individual earns more or less than \$50,000 per year. 

\vspace{2 \baselineskip}
\section{Description of our work}
We will develop a binary classification model to predict income levels. To do so, we will follow this standard methodology: 

\begin{itemize}
    \item \textbf{Data Loading \& Inspection}: Cleaning, treatment of missing values, exploration of class imbalance, encoding of categorical variables.

    \item \textbf{Exploratory Data Analysis (EDA)}: 
This is the investigation phase. We will dive deep into the data to uncover trends, visualize distributions, and analyze the relationships between different features and the target variable \emph{income}.

    \item \textbf{Feature Engineering (FE)}:
Based directly on the insights from our EDA, we will clean, transform, and create new features to prepare a dataset that our models can understand and learn from effectively.

    \item \textbf{Models training \& evaluation}: 
We will build several classification models (from a simple baseline to more complex ones), train them on our engineered data, in order to learn the relationship between features and income class. Finally, we will evaluate the performance of each model using relevant metrics to select the best one.

    \item \textbf{Models comparison}: 
Out of all prediction models, we will select the most efficient one, based on the previously calculated metrics. 
\end{itemize}

\vspace{1 \baselineskip}
Due to the imbalanced dataset the accuracy is not a relevant metric that's why we choose

\vspace{2 \baselineskip}
\section{The Dataset}
This dataset comes from the UCI Machine Learning Repository, and is available online at : 

\underline{https://www.kaggle.com/datasets/wenruliu/adult-income-dataset}

\vspace{1 \baselineskip}
The dataset contains 48,842 rows and 15 columns (14 features and 1 target variable). The features are demographic and employment-related variables. They include both numerical and categorical variables. Finally, the target variable is the income (either '$\le$50k' or '$>$50k'), that we will try to predict. 

\vspace{2 \baselineskip}
